\documentclass[12pt]{article}
\usepackage[utf8]{inputenc}
\usepackage[T1]{fontenc}
\usepackage{hyperref}
\title{Khetrani Language Resources for K-12 Curriculum Development}
\author{Compiled by Claude Sonnet 5.5 (claude-sonnet-5-5) for Andr\'as Kornai}
\date{}
\begin{document}
\maketitle

\section{Introduction}
Khetrani (Khetranki; ISO 639-3 xhe, Glottocode khet1238) is an Indo-Aryan language of north-eastern Balochistan, Pakistan. It is spoken mainly by the Khetran tribe in Barkhan District and parts of neighbouring Kohlu and Musakhel districts. It is usually placed in the Lahnda or Sindhi-Lahnda group, and Glottolog lists it under Sindhic. Reported speaker numbers conflict widely: about 15,000 (Joshua Project), about 20,000 (Wikipedia's Lahnda page), over 100,000 (2017) and 200,000 (2023) in different Wikipedia versions. Very little teaching material exists. I found no primers, textbooks, school programmes or dictionaries for Khetrani.

\section{Primary Educational Resources}
None found.

\section{Linguistic and Reference Materials}
\begin{itemize}
\item Birmani, Ali H.\ and Fasih Ahmed (2017), ``Language of the Khetrans of Barkhan of Pakistani Balochistan: A preliminary description'', \emph{Lingua} 191--192, pp.\ 3--21 (DOI 10.1016/j.lingua.2016.12.003). This is the main scholarly description I found; I could not access the paper itself. Bibliographic record: \url{https://glottolog.org/resource/reference/id/553807}
\item Wikipedia, Khetrani language: \url{https://en.wikipedia.org/wiki/Khetrani_language}
\item Wikipedia, Khetran (people): \url{https://en.wikipedia.org/wiki/Khetran}
\item Glottolog classification page: \url{https://glottolog.org/resource/languoid/id/balo1266.bigmap.html}
\item ASJP word list data: \url{https://asjp.clld.org/languages/KHETRANI}
\end{itemize}

\section{Religious and Community Resources}
Joshua Project lists the Khetran Baloch in Pakistan: \url{https://joshuaproject.net/people_groups/20248}. I found no scripture translation or community language organisation for Khetrani.

\section{Digital and Audio Resources}
I found no audio, video or software resources.

\section{Curriculum Development Context}
There is no known schooling in Khetrani, and no standard orthography was found. A curriculum would first need documentation work with the Khetran community: a word list, an orthography and recorded texts. The Birmani and Ahmed description is the natural starting point. Its relation to Sindhi, Saraiki and Punjabi is disputed in the sources, and Khetrani should not be treated as a dialect of any of them. Urdu is likely the practical school language of the region.

\end{document}